\chapter{Grundlagen}\label{ch:grundlagen}

\section{Large Language Models (LLMs) und Funktionsweise}

\subsection{Definition und Grundprinzip}

\subsection{Entwicklung}

\subsection{Kontextfenstewr als zentrales Konzept?}

\subsection{Trainingsparadigma}

\subsection{Prompting / Steuerung}

\subsection{Überleitung zu Toolcalling}

\subsection{Sicherheitsbedrohungen bei LLMs}

\subsubsection{Einordnung}

\subsubsection{Kurze Taconomie den OWASP Top 10 Rahmen}

\section{Tool Calling}

\subsection{Sicherheitsbedrohungen bei Tool Calling}

\section{Chatbots Definition}

\section{Chatbot Architekturen (technischer hintergrund -> wo angriffe drauf ausgeführt werden können)}

\section{Unternehmensinterne Chatbots: Einsatzbereich und Risiken}